% (User input window)
% (Source tag) /killo-golden
% Timestamp: 2026-02-18 15:11:15 (Asia/Singapore)
%
% Rewritten using the existing notation and auditable interfaces of this manuscript, without introducing manual numbering.

\begin{conclusion}[Discriminant and modulo-\texorpdfstring{$p$}{p} degenerations of the Lee--Yang branch polynomial: branch-value collisions at \texorpdfstring{$\{3,31,37\}$}{3,31,37}]\label{con:xi-terminal-zm-leyang-branchpolynomial-discriminant-modp-collisions}
Let the Lee--Yang cubic polynomial be
\[
P_{\mathrm{LY}}(y)=256y^3+411y^2+165y+32\in\ZZ[y],
\]
and define the total branch polynomial
\(
B(y):=y(y-1)P_{\mathrm{LY}}(y)\in\ZZ[y]
\).
Then:
\begin{enumerate}
  \item The discriminant of \(P_{\mathrm{LY}}\) is
  \[
  \boxed{\ \mathrm{Disc}(P_{\mathrm{LY}})=-3^9\cdot 31^2\cdot 37\ }.
  \]
  \item The discriminant of \(B(y)\) is
  \[
  \boxed{\ \mathrm{Disc}(B)=-2^{20}\cdot 3^{15}\cdot 31^2\cdot 37\ }.
  \]
  Therefore, branch-value merging in characteristic \(p\) can occur only for \(p\in\{2,3,31,37\}\);
  for \(P_{\mathrm{LY}}\), the actual collisions occur precisely at \(p\in\{3,31,37\}\).
  \item These degenerations have fully explicit factorization types modulo primes:
  \[
  P_{\mathrm{LY}}(y)\equiv (y-1)^3 \pmod 3,\qquad
  P_{\mathrm{LY}}(y)\equiv 8\,(y-12)^2(y-6)\pmod{31},
  \]
  \[
  P_{\mathrm{LY}}(y)\equiv -3\,(y-14)(y-6)^2\pmod{37}.
  \]
  Thus in characteristic \(3\), all three Lee--Yang branch values collapse to a single point (a triple root), while in characteristics \(31\) and \(37\), double collisions occur (one simple root and one double root).
  \item The prime \(37\) is simultaneously one of the discriminant-ramification primes for the Latt\`es critical sextic in Conclusion \ref{con:xi-terminal-zm-elliptic-lattes-critical-sextic-galois}, and the unique odd bad-reduction prime of \(\Delta(E)=37\) (Conclusion \ref{con:xi-terminal-zm-elliptic-normalization}); hence the critical dynamical structure and Lee--Yang branch geometry align in arithmetic sensitivity.
\end{enumerate}
\end{conclusion}
\begin{proof}[Proof sketch]
Discriminant factorization and modulo-\(p\) factorizations are finite algebraic computations in \(\ZZ[y]\) and \((\ZZ/p\ZZ)[y]\).
\(\mathrm{Disc}(B)\) can also be recovered from \(\mathrm{Disc}(P_{\mathrm{LY}})\) and the multiplicative formula for the resultant \(\mathrm{Res}(P_{\mathrm{LY}},y(y-1))\).
Reproducibility audit script: \texttt{scripts/exp\_fold\_zm\_elliptic\_audit.py}. This completes the proof.
\end{proof}

